\section{Problem Formulation}
Consider a finite-horizon controlled process with observation $o_t$ and
frozen policy $a_t^\pi=\pi(o_t)$. We seek an applied action $a_t$ in a bounded set
$\mathcal A$ that preserves a safety margin $h(s_{t+1})\geq0$ and stays
close to $a_t^\pi$. The policy is trained before calibration and evaluation.
The deployment filter changes its actions without updating policy weights.

The ACC state is $s=(v_e,v_l,g,v_d,\phi)$, comprising ego speed, lead
speed, inter-vehicle gap, desired speed, and grade. We linearly normalize a
protocol-defined actuator interval $[F_{\min},F_{\max}]=[-6000,4000]$ N:
$F(a)=F_{\min}+(a+1)(F_{\max}-F_{\min})/2=-1000+5000a$ N for
$a\in[-1,1]$. Under a gain shift $\gamma$, the simulator applies
$F_\gamma(a)=-1000+5000\gamma a$ N. These bounds specify the experiment;
the linear normalization is the policy--actuator interface. With sampling interval
$\Delta t=0.1$ s, longitudinal prediction uses drag, rolling resistance,
grade, and a mass-dependent acceleration \cite{rajamani2012vehicle}.
The nominal model $f_0$ approximates these effects and predicts constant lead
speed. Model hypotheses $f_\theta$ vary mass, grade, and actuator gain over
a finite set $\Theta$.

We use the constant-time-headway (CTH) spacing policy
$d_{\mathrm{CTH}}=d_0+T_hv_e$, a standard ACC design
\cite{wu2020spacing,molnar2023cruise}. Our protocol fixes
$d_0=3$ m, $T_h=1.5$ s, and $v_{\max}=30$ m/s, giving
\begin{equation}
 h_{\mathrm{CTH}}(s)=\min\{g-d_0-T_hv_e,\;v_{\max}-v_e\}.
 \label{eq:margin}
\end{equation}
The margin changes with $g$ and $v_e$, while its three protocol coefficients
remain fixed. Closing-speed and braking-distance alternatives such as
time-to-conflict and responsibility-sensitive safety (RSS) represent a richer
dynamic risk model \cite{molnar2023cruise,shalevshwartz2017rss}.
For the Euclidean norm in these physical coordinates,
$h_{\mathrm{CTH}}$ is Lipschitz with
$L_h=\sqrt{1+T_h^2}=\sqrt{3.25}$.
Calibration and evaluation use the same physical coordinates; changing their
scaling requires recalibrating the score and its Lipschitz conversion.

The episode endpoint monitors the same time-headway flag $g<3+1.5v_e$;
collision, violation duration, tracking error, and comfort provide
complementary outcomes.
